\section{Poisson三次incidenceと37次元の周囲成分}
\label{sec:pair-generic-component}
\subsection{\texttt{poisson-p4}の三次incidenceとの照合}
本節は関連原稿\cite{PoissonP4Incidence}の三次延長定理と
ファイバー計算を用いる。射影四空間の全成分数やその完全分類は仮定しない。
本稿ではPoisson構造のスカラー倍を\emph{同一視しない}。
同原稿の射影化した次元36に対し、ここでの錐の次元は37である。
この区別は、有限量子パラメータを持つ代数のモジュライを考えるために必要である。

$V^*=\langle x,y,z,u,v\rangle$、
$\Omega=dx\wedge dy\wedge dz\wedge du\wedge dv$とし、
\begin{equation}\label{eq:rat113-jacobian}
 \iota_{P(F,L,M)}\Omega=dF\wedge dL\wedge dM,
 \qquad \deg F=3,\quad \deg L=\deg M=1
\end{equation}
と置く。$C_{113}^{\rm aff}$はこの全像の被約閉包である。
標準発散作用素を$\delta$と書く。
閉三形式の分解可能性から$\delta P=0$かつ$[P,P]=0$である。

\begin{proposition}[滑らかな三次からの像と四次元ファイバー]
\label{prop:poisson-cubic-ambient-image}
滑らかな三次超曲面$X=V(F)\subset\PP^4$をPoisson部分多様体にする
非零ambient Poisson構造は、必ず\eqref{eq:rat113-jacobian}型で一般階数二である。
$X$上のPoisson構造を固定した延長スキームは被約一点である。
非零ambientを固定して$X$を忘れる向きのファイバーは、
尺度を規格化した後の
\begin{equation}\label{eq:poisson-cubic-fiber}
 \{F+g(u,v):g\in\Sym^3\langle u,v\rangle,\ V(F+g)\text{ は滑らか}\}
\end{equation}
という$\A^4$の開部分である。
従って三次Poisson対の像は$C_{113}^{\rm aff}$に稠密であり、
その次元は$41-4=37$である。
\end{proposition}
\begin{proof}
$\wedge^2T_{\PP^4}\simeq\Omega^2_{\PP^4}(5)$なので、三次超曲面への
制限の核は$H^0(\PP^4,\Omega^2_{\PP^4}(2))=0$である。
Loray--Pereira--TouzetおよびMayanskiyの分類\cite{LPT,Mayanskiy}により、
$X$上の非零Poisson構造は分解可能な定数二形式に対応し、
\eqref{eq:rat113-jacobian}の制限で与えられる。
線形な制限方程式の核が零であるから、延長はスキームとしても一点である。

射影bivectorの二次斉次liftはEuler場$E$との外積だけ曖昧だが、
$\delta(E\wedge Y)=5Y-(\operatorname{tr}Y)E$より
発散零の代表は一意である。その代表の一次Casimir空間は
$U=\langle L,M\rangle$に等しい。実際、独立な第三の一次Casimirがあれば
$F$は三つの一次式だけの多項式になり、$V(F)$は錐となって滑らかさに反する。
従ってambientから$U$を回復する。同じ$U$に対する二つのJacobianが
比例する条件は$F'-cF\in\Sym^3U$である。尺度を除くと
\eqref{eq:poisson-cubic-fiber}を得る。
滑らかな三次式は係数空間の稠密開であり、
Grassmannianの6次元と、非射影化した相対Jacobianの31係数を合わせて
像の次元は37である。Poisson対の次元41は、三次式の射影係数34と
固定した三次上のPoisson錐の次元7の和でもある。
\end{proof}

\subsection{全Poissonスキーム内の接空間と一般被約性}
\begin{theorem}[\texorpdfstring{$\mathrm{Rat}(1,1,3)$}{Rat(1,1,3)}の独立な成分証明]
\label{thm:rat113-smooth-component}
$C_{113}^{\rm aff}$は、五変数の発散零二次Poisson構造の
\emph{全スキーム}の37次元既約成分であり、その一般点は滑らかである。
従って射影四空間の非射影化Poissonスキームでも同じ結論が成り立つ。
\end{theorem}
\begin{proof}
二次bivectorの係数空間は$10\cdot15=150$次元である。
全ての狭義Jacobi方程式$\frac12[P,P]=0$の350係数と、
発散$\delta P=0$の25係数を用いる。
次の整数係数三次式を取る：
\begin{align}\label{eq:rat113-test-F}
 F_*={}&3(x^3+y^3+z^3)+5xyz+ux^2+2vy^2+3(u+v)z^2\notag\\
 &+(u^2+uv+2v^2)x+(2u^2+3uv+v^2)y
 +(u^2+4uv+5v^2)z.
\end{align}
$P_*=P(F_*,u,v)$における方程式の微分を$M$とする。
列は十個の辞書式添字対と15個の二次単項式、
行は十個の辞書式添字三つ組と35個の三次単項式、
続いて五つの発散成分と五つの一次単項式の順である。
例えばJacobi成分の微分は
\begin{equation}\label{eq:rat113-jacobi-differential}
 (DQ)_{ijk}=\sum_{\mathrm{cyc}(i,j,k)}\sum_{l=1}^{5}
 \left(Q_{il}\partial_l(P_*)_{jk}
       +(P_*)_{il}\partial_lQ_{jk}\right)
\end{equation}
で与えられる。添字を$1,\ldots,5$として変数順$x,y,z,u,v$を使う。

一方、$F_*$のうち純$u,v$単項式を除いた31係数、および
$u\mapsto u+s_1x+s_2y+s_3z$、
$v\mapsto v+t_1x+t_2y+t_3z$の六係数による族の微分を$S$とする。
$M$は$375\times150$、$S$は$150\times37$の整数行列であり、
Jacobian族が恒等的にPoissonかつ発散零なので$MS=0$である。
再現コードはこの積も整数上で直接確認する。

\nolinkurl{computations/verify_rat113_tangent.py}は両行列を元の公式から生成し、
\nolinkurl{computations/rat113_tangent_certificate.json}の行・列添字で
指定した小行列式について
\begin{equation}\label{eq:rat113-minor-certificates}
 \det M_{113}\equiv844194,\qquad
 \det S_{37}\equiv948131\pmod{1000003}
\end{equation}
を得る。これらは整数行列式の\emph{非零}証明である。
従って$\rank_{\mathbb Q}M\ge113$、$\rank_{\mathbb Q}S=37$であり、
整数上の恒等式$MS=0$と合わせて$\rank_{\mathbb Q}M=113$となる。
有限体での零から特性零の恒等式を推論してはいない。
よって全スキームの接空間は$150-113=37$次元である。

この点を通るJacobian族の像は、$S$の階数37から局所次元37を持つ。
したがって全スキームの局所次元と接空間次元が一致し、この点は正則点である。
Jacobian源は既約であるから、その37次元像の閉包がこの点を通る
唯一の既約成分となり、滑らかな開部分を持つ。

最後に、発散零二次bivectorの空間は
$H^0(\PP^4,\wedge^2T_{\PP^4})$に線形同型である。
発散零の$P$に対して射影Jacobiが零なら、
$[P,P]=E\wedge A$と$\delta[P,P]=0$を用いて
$0=5A-E\wedge\delta A$、従って$E\wedge A=0$を得る。
これは係数の任意の基底環上で成立するため、Jacobiイデアルも一致する。
従って上記は射影Poissonスキームそのものについての証明である。
\end{proof}

\subsection{規格化CY変形の局所成分への移行}
Pym\cite[Definition 2.8, Proposition 2.9, Theorem 2.10]{Pym}の
$\operatorname{QDef}_0(k^5)$は、\emph{可換多項式環の点における}
Koszul Calabi--Yau変形の芽である。
同文献第2節のこれらの結果は任意の$\dim V$について述べられており、
第3節の四生成元の六成分分類とは区別する。
以下の成分性は、この局所変形空間に関する主張である。

$V=k^5$とし、$V^{\otimes2}=\wedge^2V\oplus\Sym^2V$で規格化して
\[
 R_\Pi=\langle a\otimes b-b\otimes a-\Symm(\Pi(a,b)):a,b\in V\rangle,
 \qquad B_\Pi=T(V)/(R_\Pi)
\]
と置く。
\begin{proposition}[graph成分と一般被約性]\label{prop:quantum-113-component}
$\Pi\mapsto R_\Pi$は$C_{113}^{\rm aff}$を二次関係Grassmannianの
標準アフィンchartへ閉埋込みする。その像$Q_{113}$は37次元で、
可換点の規格化CY変形の芽において対応する被約既約成分をなす。
この成分の一般点における全局所環は被約かつ正則である。
これは全AS正則$\mathbb Z$-代数の大域モジュライの成分性を意味しない。
\end{proposition}
\begin{proof}
$R_\Pi$は$\wedge^2V\to\Sym^2V$のgraphなので$\Pi$を一意に回復でき、
標準liftが線形であることから閉埋込みとなる。
最高次平面三次制限が滑らかな相対Jacobian点は稠密であり、
その点では補題\ref{lem:pbw}と\ref{lem:CY}がKoszul CY代数を与える。
CY変形の条件はこの変形芽では閉条件なので、閉包の芽も含まれる。
それを含む被約既約成分の芽を$D$とする。
一次項が恒等な形式比較の下で$D$の接錐の被約台はあるPoisson既約成分$C_D$である。
graph像は斉次錐なので
$C_{113}^{\rm aff}=TC_0Q_{113}\subset TC_0D=C_D$となる。
定理\ref{thm:rat113-smooth-component}の成分の極大性から両者は等しく、次元も37である。
従って$D$と$Q_{113}$の被約な既約芽は等しい。
Pymの比較は被約台だけの比較ではなく、Poissonスキームと規格化CY変形の
形式的、従ってある解析的な芽の同型である。
定理\ref{thm:rat113-smooth-component}による一般正則性も対応する成分へ移る。
一般点では他の成分との交叉を除けるので、全局所環も正則となる。
これはKontsevich級数自体がgraph式と一致するという主張ではない。
\end{proof}

\begin{lemma}[相対Jacobian括弧の全Poisson中心]
$R=k[u,v]$とし、$F\in R[x,y,z]$の最高active次数部分$f_3$が既約とする。
$R$を中心、$\{x,y\}=F_z$などと定めると$Z_{\rm Pois}(R[x,y,z])=R[F]$である。
\end{lemma}
\begin{proof}
中心元の最高次部分$g_m$は$df_3\wedge dg_m=0$を満たす。
Euler場との縮約から$3f_3dg_m=mg_mdf_3$。
$f_3$は素元で、少なくとも一つの非零偏微分を割らないので$f_3\mid g_m$である。
これを繰り返すと$m=3j$、$g_m=cf_3^j$ ($c\in R$)となる。
$g-cF^j$へ次数帰納法を適用すればよい。逆包含はCasimir性による。
\end{proof}
\begin{theorem}[量子成分の関数体上の中心]
\label{thm:pair-generic-center}
$K=k(Q_{113})$とし、一般ファイバーを$B_\eta$とすると
\[
 Z(B_\eta)=K[a,b,c],\qquad\deg a=\deg b=1,\quad\deg c=3.
\]
$B_\eta$は整域で、中心の次数1,2,3の次元は2,3,5となる。
全次数の中心が一つのZariski開部分で一定であるとは主張しない。
\end{theorem}
\begin{proof}
中心変数を厳密に保つ相対Jacobian量子化と中心元$G_\hbar\equiv F\pmod\hbar$を取る。
次数$d$の形式中心元を$\hbar$で還元すると前補題のPoisson中心に入り、
$u,v,G_\hbar$の同じ重みの多項式を引けば差は$\hbar$で割れる。
反復と完備性から
\[
 Z(B_\hbar)_d=\bigoplus_{i+j+3l=d}k[[\hbar]]u^iv^jG_\hbar^l.
\]
独立性は関係式の最小$\hbar$次数を取ることで従う。

この計算を成分の一般点へ戻す。
標準次数付き整域$A$の頂点完備化から、支配的パラメータ表示$T\to\Spec A$を用いて
\[
 \widehat A\longrightarrow k(T)[[\hbar]],\quad
 \sum a_n\longmapsto\sum\hbar^n\phi^*a_n
\]
を作ると、最初の非零斉次係数が消えないので単射である。
これをPoisson錐と形式量子化比較に適用し、$K$を含む体
$L=k(T)((\hbar))$を得る。$T$はpencilの6座標と三次式の31係数で取れる。
各次数の中心は有限自由空間の交換子行列の核なので体拡大と可換し、
$Z(B_\eta)_d\otimes_K L=L[u,v,G_\hbar]_d$となる。
$Z(B_\eta)_1$の基底$a,b$を取り、$c\in Z(B_\eta)_3\setminus\Sym^3\langle a,b\rangle$を選ぶと、
$L$上では$c=\lambda G_\hbar+r(u,v)$ ($\lambda\ne0$)である。
従って$K[a,b,c]\to Z(B_\eta)$は忠実平坦な体拡大後に同型、ゆえに元から同型である。
整域性も形式PBW量子化への埋込みから従う。
\end{proof}
\begin{corollary}[支配的incidenceと有限次数の良い開部分]
$Q_{113}$の次数三中心直線のincidenceには一般ファイバー$\PP^4_K$を持つ
41次元の支配的既約成分$I_{113}$が一意に存在する。
中心の次数1,3とHilbert関数の次数4以下が一般値を取る開部分では、
これはrank5の中心核束の$\PP^4$束である。
一般の有理商の次元は周囲で13、対で17、忘却の一般ファイバーは4次元である。
\end{corollary}
\begin{proof}
中心の有限行列の非零minorを反転するとrank5の核束を得る。
その射影束は既約で一般ファイバー$\PP^4_K$の閉包を与え、支配的成分は他にない。
相対Jacobianパラメータ族の6+35=41次元の像がこれを支配する。
六分岐値を持つ対角族の周囲は射影stabilizerが有限であることを
第\ref{sec:pair-diagonal-classification}節で証明する。
従って$\PGL_5$の一般軌道は24次元である。
Rosenlichtの定理により37−24=13、41−24=17を得て、忘却の一般ファイバーは4次元となる。
この計算は商多様体の有理性や一般stabilizerの自明性を含まない。
\end{proof}
